\chapter*{\centering Résumé}
\addcontentsline{toc}{chapter}{RÉSUMÉ}

\hs La souscription d'assurance maladie dans les marchés émergents tels que le Cambodge se heurte à des défis structurels, notamment un faible taux de pénétration de l'assurance, une infrastructure de santé fragmentée et des dossiers de santé longitudinaux limités. Les systèmes de souscription statiques fondés sur des règles sont sous-optimaux, car ils ignorent les interactions entre variables, ne peuvent s'adapter à la dérive du portefeuille et risquent de discriminer involontairement certains segments démographiques. Ce mémoire examine si les bandits contextuels --- une classe d'algorithmes d'apprentissage en ligne pour la prise de décision séquentielle en situation d'incertitude --- peuvent remplacer les règles statiques dans la souscription d'assurance maladie des marchés émergents.

Cette recherche conçoit et met en {\oe}uvre un cadre de bandits contextuels reposant sur trois algorithmes (LinUCB, LinTS et Epsilon-Greedy) et les compare à une référence statique fondée sur des règles XGBoost, sur un jeu de données synthétique de 2{,}000 candidats cambodgiens à l'assurance maladie, calibré sur l'Enquête démographique et de santé du Cambodge 2021--22 \citep{NIS2023}. Un simulateur de récompense actuarielle fondé sur le profit évalue quatre décisions de souscription (standard, surprime, refus, renvoi)~; des garde-fous fondés sur l'Indice de Stabilité de la Population (PSI) à fenêtre glissante et la règle des quatre cinquièmes de l'EEOC américaine surveillent conjointement l'équité régionale et professionnelle. Tous les résultats principaux sont présentés sous forme de moyennes sur 20 graines indépendantes, accompagnées de tests des rangs signés de Wilcoxon appariés et d'intervalles de confiance bootstrap à 95\%.

Quatre expériences contrôlées valident le cadre proposé. EXP-005 montre que LinUCB atteint une récompense cumulée de \$90{,}540 sur 5{,}000 tours, contre \$72{,}292 pour la référence statique ($+25\%$, Wilcoxon $p < 0.001$, $d$ de Cohen $= 2.98$), et réduit le regret moyen des 500 derniers tours de \$9.01 à \$2.20 par tour. EXP-006 confirme que le PSI maximal en fenêtre glissante demeure dans la bande VERTE/AMBRE pour la région (0.082) comme pour la profession (0.123), tandis que la règle des quatre cinquièmes de l'EEOC est respectée, avec des ratios de parité de 85.7\% et 90.1\% respectivement. EXP-007 classe les algorithmes selon le regret cumulé~: LinTS et LinUCB sont statistiquement indiscernables et surpassent tous deux nettement les références statique et à exploration uniforme. EXP-008 enveloppe LinUCB dans une couche de supervision humaine (\emph{human-in-the-loop}) et augmente la récompense cumulée de 14.8\% par rapport à la référence REFER mathématique, pour un coût de révision humaine de 2.2\% de la récompense. Une analyse de robustesse a posteriori identifie un plafond de performance théorique~: une politique constante inadmissible (AlwaysRATED) qui surpasse les deux bandits de 30 à 33\% --- ce qui borne ce résultat, sans toutefois le renverser.

Les résultats fournissent une preuve de concept rigoureuse~: les bandits contextuels peuvent surpasser la souscription statique fondée sur des règles qu'ils remplaceraient sur le plan de la rentabilité tout en préservant l'équité démographique --- et demeurer d'un coût de calcul négligeable, le modèle sous-jacent étant un estimateur linéaire à 34 variables mis à jour analytiquement~; le plafond constant inadmissible évoqué ci-dessus borne cette contribution face aux alternatives admissibles, sans toutefois l'invalider. Le mémoire apporte un banc d'essai expérimental reproductible à 20 graines, un jeu de données synthétique calibré sur le Cambodge, une surveillance de l'équité par PSI à fenêtre glissante et règle des quatre cinquièmes de l'EEOC, ainsi qu'un dispositif de supervision humaine, susceptibles d'éclairer les recherches futures et les pratiques de l'industrie en matière d'assurance algorithmique sur les marchés émergents.

\bigskip
\noindent\textbf{Mots-clés :} Bandits contextuels, LinUCB, échantillonnage de Thompson, apprentissage par renforcement, souscription d'assurance maladie, Cambodge, équité algorithmique, Indice de Stabilité de la Population, science actuarielle, apprentissage en ligne
